\section{The linear-probe backbone-selection diagnostic}
\label{app:probe}

This appendix details the linear-probe step referred to in \cref{sec:pretrain} for choosing among
backbone candidates (the SSL-pretrained ResNets, the supervised ResNet-50 baseline) and among BN-only
recalibration variants, before any M1 machinery is trained. It is relevant specifically to
\nameref{sec:archB} (\cref{sec:archB:gap}): \nameref{sec:archA} does not introduce a separate, frozen
$\tilde x=\mathrm{BB}(x)$ in the first place, so there is nothing here for it to check. The question
this probe answers is narrow, deliberately not causal, and -- this is worth stating plainly given
\cref{sec:archB:gap} -- deliberately not about the property \cref{thm:ng} actually needs from
$\mathrm{BB}\circ g$: given a frozen, cached embedding $\tilde x$, is content information linearly
readable, and does location/style information survive at all. Linear readability is a much weaker
property than the injectivity \cref{sec:archB:gap} identifies as what the identifiability argument
requires; a backbone can pass both probes below while still collapsing many distinct $(Z_c,Z_s)$ onto
embeddings a linear classifier can separate but a decoder cannot invert. This appendix is therefore a
feature-extraction sanity check that gates whether \nameref{sec:archB}'s heuristic is worth attempting
for a given backbone at all, not a substitute for the M1 estimator itself, and not evidence that the
heuristic's output, if it converges, bears any proven relationship to the true $(Z_c,Z_s,\Gc)$.

\subsection{Why a linear probe, and why two of them}
A linear probe -- fitting a plain linear classifier on top of a frozen embedding, with no
backpropagation into the backbone -- tests what is already linearly present in the features, rather
than what a sufficiently powerful downstream model could eventually extract from them. This matches
the concern raised about \nameref{sec:archB} in \cref{sec:archB:gap}: whatever is not present in 
the cached embedding is a hard ceiling that no downstream causal machinery can recover
(\cref{sec:archB}, cons of \nameref{sec:archB}). The probe is run once per backbone/BN-variant 
candidate, on cached features only, before committing to \nameref{sec:archB}'s heuristic for that
candidate.

Two separate probes are needed because content and style are different questions:
\begin{itemize}[leftmargin=1.5em]
\item \textbf{Species accuracy.} Input: the cached embedding of an image (layer3 pooled concatenated
      with the final pooled layer, per \cref{sec:archB}). Label: species (10 classes). This checks that
      content information survived extraction at all; the target is high accuracy.
\item \textbf{Within-species location accuracy.} Input: the same cached embedding, restricted to images
      of one species at a time (i.e.\ conditioned on species). Label: training location
      (L38/L43/L46). Species and location are correlated in the raw data, so an unconditional
      location probe could score well purely by piggybacking on species and would say nothing about
      style; conditioning on species asks the real question, whether the embedding still separates
      locations once species is held fixed. The target here is not high accuracy but accuracy clearly
      above chance -- chance-level would mean the backbone, even after BN-only recalibration, erased
      the location/style cues the style block $Z_s$ needs.
\end{itemize}

\subsection{Data split}
\label{app:probe:split}
\begin{itemize}[leftmargin=1.5em]
\item \textbf{L100 is never used by either probe}, for fitting or for scoring. This diagnostic selects
      a backbone using only train-location data; L100 stays reserved for evaluating the eventual M1
      pipeline, consistent with its role everywhere else in this note.
\item \textbf{Within L38/L43/L46, use stratified $k$-fold cross-validation}
	\footnote{``Stratified jointly by species and location" means the folds are built to preserve in 
	each fold the proportions of each species$\times$location combination in the full dataset, so every 
	fold has a representative slice of, say, ``raccoon at L43" and not just ``raccoon" in general. That 
	matters specifically for the within-species location probe: since it's scored per species, you need 
	each species to have reasonable location coverage in every fold, or some folds simply can't produce 
	a meaningful location accuracy for the rarer species.}
	(e.g.\ $k=5$), stratified
      jointly by species and location, rather than a single train/test split. Two reasons a single
      split is not enough here: species classes are imbalanced (raccoon dominant; bird/dog/squirrel
      data-limited), so one random split risks leaving too few examples of the rare classes in the
      held-out fold; and the within-species location probe is evaluated per species, so each
      species$\times$location cell is smaller than the marginal class counts and some cells are already
      sparse. Cross-validation uses all train-location data for both fitting and scoring across folds
      and gives a less noisy estimate for comparing candidates than a single split would.
\item In every fold, the linear classifier is fit on the training portion of that fold only and scored
      on the held-out portion; no image contributes to both fitting and scoring for the same fold.
\end{itemize}

\subsection{Metric}
Report \textbf{macro accuracy} (per-class accuracy, averaged across classes) for both probes, not raw
accuracy, for the same reason macro accuracy is the chosen metric for the main species classifier
(\cref{sec:setting}): with raccoon dominant and bird/dog/squirrel data-limited, raw accuracy would be
dominated by the majority class and could mask whether rarer species, or rarer
species$\times$location cells, are actually separable in the embedding. For the within-species probe,
compute macro accuracy over locations within each species, then average over species (weighting
species equally, not by sample count), so that a species with few images does not vanish from the
result.

\subsection{Procedure}
\begin{enumerate}[label=(\arabic*),leftmargin=2em]
\item For each candidate (backbone $\times$ BN-recalibration variant), take the cached embeddings for
      all train-location images (\cref{sec:archB}), excluding L100.
\item Run stratified $k$-fold cross-validation (\cref{app:probe:split}) with a linear classifier
      (e.g.\ multinomial logistic regression) to get:
      \begin{itemize}[leftmargin=1.3em]
      \item macro species accuracy, label $=$ species, over all train-location images;
      \item macro within-species location accuracy, label $=$ location, computed separately within
            each species and averaged over species.
      \end{itemize}
\item Record both numbers per candidate.
\end{enumerate}

\subsection{Decision rule}
Prefer the candidate with high macro species accuracy and macro within-species location accuracy
clearly above chance (chance being $1/|\Lset_t|$ within each species' available locations, averaged
the same way as the metric). A candidate that is high on species accuracy but at chance on location
accuracy is discarded regardless of species accuracy: per the feature-extraction ceiling argument of
\cref{sec:archB}, no amount of M1 estimator design can recover style information that the backbone
never retained. Among candidates that pass both checks, this diagnostic does not by itself rank them
further; \cref{sec:pretrain} notes there is no first-principles reason to prefer one self-supervised
pretraining method over another for style preservation, so ties are broken by whichever gives the best
combined result empirically. In every case, a pass narrows candidates to those that plausibly retain
enough signal for \nameref{sec:archB}'s heuristic to have a chance; per \cref{sec:archB:gap}, it does
not certify that any candidate yields a $\mathrm{BB}\circ g$ close enough to injective for
\cref{thm:ng} to apply, since no test proposed here checks that property directly.
